\PassOptionsToPackage{hyphens}{url} 
\usepackage{hyperref}
% \usepackage{cite}
\usepackage{amsmath,amssymb,amsfonts}
\newcommand{\proves}{\vdash}
\newcommand{\adv}{\mathcal{A}}
\usepackage{textcomp}
\usepackage[utf8]{inputenc}
\usepackage{xcolor}
\usepackage{lscape}
\usepackage{pdflscape}
% \usepackage[pdftex]{graphicx}
\usepackage{graphicx}
\usepackage{booktabs}
\usepackage{dirtytalk}
\usepackage{longtable}
\usepackage{wasysym}
\usepackage{multirow}
\usepackage[inline]{enumitem}
\usepackage[utf8]{inputenc}
\usepackage{csquotes}
% \usepackage{acronym}
\usepackage{tikz}
\usetikzlibrary{turtle}
\usepackage{acro}
\usepackage[english]{babel}
\usepackage{tabularx}
\usepackage{longtable}
\usepackage{colortbl}
% \usepackage{filecontents}
\usepackage{array}
\usepackage{makecell}
\usepackage[normalem]{ulem}
\useunder{\uline}{\ul}{}
% \usepackage[ruled,vlined]{algorithm2e}
\usepackage{colortbl}

\usepackage{algorithm}
\usepackage{algpseudocode}
\usepackage{xcolor}

\usepackage{pifont}


\usepackage{placeins}

\usepackage[numbers]{natbib}

\usepackage{rotating}

\usepackage{xurl}

\usepackage{etoolbox}

\usepackage{needspace}

% \usepackage[acronym]{glossaries}
\usepackage[acronym]{glossaries-extra}
\makeglossaries



% \usepackage[numbers]{natbib}

% \usepackage{url}
% \usepackage{usenix2019_v3}
% \usepackage{usenix,epsfig}

% Reduce vertical spacing between acronym entries


% \setglossarystyle{list} % clean list style (or try 'long' or 'super')

% \setlength{\glsdescwidth}{0.75\linewidth} % optional: align definitions
% \setlength{\glsentrysep}{0.4em}


% \renewenvironment{theglossary}%
%   {\begin{description}[style=unboxed,leftmargin=3em,labelwidth=4em]}%
%   {\end{description}}

% % Customize item formatting (optional)
% \renewcommand*{\glossaryentryfield}[5]{%
%   \item[\glsentryshort{##1}] ##2%
% }

\def\tsc#1{\csdef{#1}{\textsc{\lowercase{#1}}\xspace}}
\tsc{WGM}
\tsc{QE}
\tsc{EP}
\tsc{PMS}
\tsc{BEC}
\tsc{DE}

\setlength{\parindent}{0pt} % Set paragraph indentation to zero

\NewDocumentCommand{\rot}{O{45} O{1em} m}{\makebox[#2][l]{\rotatebox{#1}{#3}}}

\newtheorem{definition}{Definition}[section]

% Look into this later for having items automatically bold
% \renewcommand{\item}{\bfseries\origitem}

\newcommand{\oldstuff}[1]{}
\newcommand{\old}[1]{}


% empty half and full circles
\newcommand*\emptycirc[1][1ex]{\tikz\draw (0,0) circle (#1);} 
\newcommand*\halfcirc[1][1ex]{%
  \begin{tikzpicture}
  \draw[fill] (0,0)-- (90:#1) arc (90:270:#1) -- cycle ;
  \draw (0,0) circle (#1);
  \end{tikzpicture}}
\newcommand*\fullcirc[1][1ex]{\tikz\fill (0,0) circle (#1);} 


% create new “smaller” versions:
\newcommand{\smallfullcirc}{\raisebox{-0.1em}{\scalebox{0.75}{\fullcirc}}}
\newcommand{\smallemptycirc}{\raisebox{-0.1em}{\scalebox{0.75}{\emptycirc}}}
\newcommand{\smallhalfcirc}{\raisebox{-0.1em}{\scalebox{0.75}{\halfcirc}}}


% test for colours 
\newcommand{\FilledCircle}{\textcolor{green}{\ding{108}}} % filled circle in green

% half circles in tables
\newcommand*\lefthalfcircle[1][1ex]{%
  \begin{tikzpicture}
    \draw[fill] (0,0) arc[start angle=180,end angle=360,radius=#1] -- cycle;
  \end{tikzpicture}}

\newcommand*\righthalfcircle[1][1ex]{%
  \begin{tikzpicture}
    \draw[fill] (0,0) arc[start angle=0,end angle=180,radius=#1] -- cycle;
  \end{tikzpicture}}

  \newcommand*\tophalfcircle[1][1ex]{%
  \begin{tikzpicture}
    \draw[fill] (0,0) arc[start angle=90,end angle=270,radius=#1] -- cycle;
  \end{tikzpicture}}

\newcommand*\bottomhalfcircle[1][1ex]{%
  \begin{tikzpicture}
    \draw[fill] (0,0) arc[start angle=270,end angle=450,radius=#1] -- cycle;
  \end{tikzpicture}}
  % -----


  % -=-=-=-=-=-=-

% heatmap colors
\definecolor{r0}{rgb}{1, 1.00, 1.00}
\definecolor{r1}{rgb}{1, 0.90, 0.90}
\definecolor{r2}{rgb}{1, 0.80, 0.80}
\definecolor{r3}{rgb}{1, 0.70, 0.70}
\definecolor{r4}{rgb}{1, 0.60, 0.60}
\definecolor{r5}{rgb}{1, 0.50, 0.50}
\definecolor{r6}{rgb}{1, 0.40, 0.40}
\definecolor{r7}{rgb}{1, 0.30, 0.30}
\definecolor{r8}{rgb}{1, 0.20, 0.20}
\definecolor{r9}{rgb}{1, 0.10, 0.10}
\definecolor{r10}{rgb}{1, 0.00, 0.00}
\definecolor{g0}{rgb}{1.00, 1, 1.00}
\definecolor{g1}{rgb}{0.90, 1, 0.90}
\definecolor{g2}{rgb}{0.80, 1, 0.80}
\definecolor{g3}{rgb}{0.70, 1, 0.70}
\definecolor{g4}{rgb}{0.60, 1, 0.60}
\definecolor{g5}{rgb}{0.50, 1, 0.50}
\definecolor{g6}{rgb}{0.40, 1, 0.40}
\definecolor{g7}{rgb}{0.30, 1, 0.30}
\definecolor{g8}{rgb}{0.20, 1, 0.20}
\definecolor{g9}{rgb}{0.10, 1, 0.10}
\definecolor{g10}{rgb}{0.00, 1, 0.00}

% Orange gradient colors
\definecolor{o0}{rgb}{1.00, 0.90, 0.70}  % Lightest orange
\definecolor{o1}{rgb}{1.00, 0.85, 0.55}
\definecolor{o2}{rgb}{1.00, 0.75, 0.40}
\definecolor{o3}{rgb}{1.00, 0.65, 0.25}
\definecolor{o4}{rgb}{1.00, 0.55, 0.10}
\definecolor{o5}{rgb}{1.00, 0.50, 0.00}  % Standard orange
\definecolor{o6}{rgb}{0.90, 0.45, 0.00}
\definecolor{o7}{rgb}{0.80, 0.40, 0.00}
\definecolor{o8}{rgb}{0.70, 0.35, 0.00}
\definecolor{o9}{rgb}{0.60, 0.30, 0.00}
\definecolor{o10}{rgb}{0.50, 0.25, 0.00} % Darkest orange


  % -=-=-=-=-=-=-


% special A font in latex
\newcommand{\specialA}{%
    \begin{tikzpicture}[baseline=(A.base)]
        \node (A) at (0,0) {\fontsize{12pt}{12pt}\selectfont A};
        \draw (A.base west) -- (A.base east);
        \draw (A.north west) -- (A.north east);
        \draw (A.north east) -- (A.base east);
        \draw (A.north west) -- (A.base west);
    \end{tikzpicture}%
}
  % -----


\newcommand{\ynote}[1]{{$\langle${\textcolor{red}{YF: \textbf{#1}}}$\rangle$}}
\newcommand{\yathin}[1]{{$\langle${\textcolor{blue}{YP: \textbf{#1}}}$\rangle$}}

\newcommand{\quotes}[1]{``#1''}
\newcommand{\paragraphtitle}[1]{\paragraph{#1}}

\newcommand{\orange}[1]{\textcolor{orange}{#1}}
\newcommand{\red}[1]{\textcolor{red}{#1}}
\newcommand{\teal}[1]{\textcolor{teal}{#1}}
\newcommand{\olive}[1]{\textcolor{olive}{#1}}
\newcommand{\gray}[1]{\textcolor{gray}{#1}}
\newcommand{\darkgray}[1]{\textcolor{darkgray}{#1}}
\newcommand{\lightgray}[1]{\textcolor{lightgray}{#1}}

\makeatletter
\let\disable\@secondoftwo
\addto\extrasenglish{%
  \renewcommand{\sectionautorefname}{\S\@gobble}%
}	
\addto\extrasenglish{%
  \renewcommand{\subsectionautorefname}{\S\@gobble}%
}
\addto\extrasenglish{%
  \renewcommand{\subsubsectionautorefname}{\S\@gobble}%
}
\addto\extrasenglish{%
  \renewcommand{\paragraphautorefname}{\S\@gobble}%
}
\makeatother

% The preceding line is only needed to identify funding in the first footnote. If that is unneeded, please comment it out.

\AtBeginDocument{%
  \providecommand\BibTeX{{%
    \normalfont B\kern-0.5em{\scshape i\kern-0.25em b}\kern-0.8em\TeX}}}



    % ---


% \usepackage{colortbl}
% \usepackage{amsmath,amsfonts}
% \usepackage{array}
% \usepackage{algorithmic}
% \usepackage{textcomp}
% \usepackage{xcolor}
% \usepackage{lscape}
% \usepackage{pdflscape}
% \usepackage{graphicx}
% \usepackage{booktabs}
% \usepackage{dirtytalk}
% \usepackage{longtable}
% \usepackage{wasysym}
% \usepackage{multirow}
% \usepackage{enumitem}
% \usepackage{booktabs}
% \usepackage{hyperref}
% \usepackage{tikz}
% \usepackage{tabularx}
% \usepackage{url}
% \usepackage{microtype}
% \usepackage{diagbox}

% \usetikzlibrary{turtle}